% Laget automatisk av elevark/lag_elevark.py fra sim/avboyning.html. Endringer her blir overskrevet.
\documentclass[11pt]{article}
\usepackage{elevark}
\begin{document}
\elevarkhode{Fysikk 2 · Elektriske felt og krefter}{Akselerasjon og avbøyning}{https://fysikklab.no/sim/avboyning.html}{Simuleringen}

En ladd partikkel starter i ro og blir akselerert gjennom spenningen $U_{\mathrm{a}}$. Så kommer den inn midt mellom to plater med spenningen $U_{\mathrm{d}}$, der feltet står vinkelrett på farten. Banen mellom platene ligner et vannrett kast, men det er den elektriske kraften som bøyer den av.

\section*{Oppgaver}

Start med et proton, $U_{\mathrm{a}} = 400\,\mathrm{V}$, $U_{\mathrm{d}} = 10\,\mathrm{V}$, $d = 4{,}0\,\mathrm{cm}$ og $l = 20\,\mathrm{cm}$. $e = 1{,}60 \cdot 10^{-19}\,\mathrm{C}$, $m_{\mathrm{p}} = 1{,}67 \cdot 10^{-27}\,\mathrm{kg}$ og $m_{\mathrm{e}} = 9{,}11 \cdot 10^{-31}\,\mathrm{kg}$.

\begin{oppgaveliste}
\oppgave \fase{test}{Regn}Bruk energibevaring, $qU_{\mathrm{a}} = \tfrac{1}{2}mv_{0}^{2}$, til å regne ut farten protonet har når det kommer inn mellom platene. Sjekk med simuleringen.
\svar
\oppgave \fase{test}{Regn}Regn ut feltstyrken, kraften og akselerasjonen mellom avbøyningsplatene. Hvor lang tid bruker protonet gjennom platene, og hvor langt er det bøyd av når det kommer ut?
\svar
\oppgave \fase{gjett}{Gjett}Kjør protonet, og bytt så til et elektron med de samme spenningene. Elektronet er nesten 2000 ganger lettere. Hvor mye mer tror du det blir bøyd av? Test, og prøv også en alfapartikkel.
\svar
\oppgave \fase{forklar}{Forklar}Vis med formler at avbøyningen ut av platene er $y = U_{\mathrm{d}} l^{2} / (4 d U_{\mathrm{a}})$. Bruk dette til å forklare resultatet i oppgave 3.
\svar
\oppgave \fase{test}{Regn}Sett lengden til 50~cm. Hvor langt inn mellom platene kommer protonet før det treffer den negative plata? Regn først, og test så.
\svar
\oppgave \fase{test}{Undersøk}Sett lengden tilbake til 20~cm. Hvor stor må $U_{\mathrm{d}}$ være for at partikkelen akkurat skal treffe enden av plata? Hvordan endres svaret hvis du halverer $U_{\mathrm{a}}$?
\svar
\oppgave \fase{utfordring}{Utfordring}I simuleringen er tyngden neglisjert. Sammenlign tyngden og den elektriske kraften på protonet mellom platene. Var det rimelig å se bort fra tyngden?
\svar
\end{oppgaveliste}

\end{document}
